\section{Arithmetic symmetric frontier cost and determinant growth}
\label{FD:FD:section}

This section proves finite restrictions on the spectral positions of the
original arithmetic packet.  It retains the actual theta source, the measure
\(\lvert 2\xi/h\rvert^2dy/(2\pi)\), all zero multiplicities, and the unscaled
orbit coordinates.  No purity assertion is used.
Stars denote conjugate transpose in the displayed finite coefficient
bases, and Hilbert adjoint for maps into the source Hilbert space.
For Hermitian matrices, \(U\preceq V\) means that \(V-U\) is positive
semidefinite.  The one-factor quotient coefficient basis is always
\([1]_h,[t]_h,\ldots,[t^{d-1}]_h\), with its specified tensor and orbit
inclusions below.

\paragraph{Authoritative source pins.}
The retained construction is in the original messages A1829 and A1859 of
\texttt{audit\_segment\_U0041\_U0054.md}.  The line numbers here are literal
lines in that source, whose SHA-256 is recorded in the accompanying audit.
A1829, lines 4956--5071, specifies the source, the arithmetic function,
the full-jet map, and its exact polynomial kernel; lines 5073--5169 give
the canonical minimum and its metric; lines 5383--5561 give the joint
tensor source and actual next relation layer.  A1859, lines 5693--5874,
gives the unscaled polynomial coordinates, the relation isomorphism,
and the exact metric update; lines 5876--6040 give the recurrence and
finite lowering bound; lines 6042--6197 give the signed symmetric
summand and its unscaled orbit factors.  The aggregate inequalities
proved below are new consequences of those retained objects.

\subsection{The retained arithmetic source and its canonical finite metric}

Keep the original operator, theta map, and arithmetic function
\begin{equation}
\begin{gathered}
D=-x\partial_x,\qquad
\Theta\phi(x)=\sum_{m\ne0}\phi(mx),\\
\phi_*(x)=(4\pi^2x^4-6\pi x^2)e^{-\pi x^2},\qquad
f_0=\Theta\phi_*,\\
g(s)=\mathcal Mf_0(s)=2\xi(s)
=s(s-1)\pi^{-s/2}\Gamma(s/2)\zeta(s).
\end{gathered}
\tag{FD.1}\label{FD:FD:arithmetic}
\end{equation}
Here \(\mathcal MF(s)=\int_0^\infty F(x)x^{s-1}\,dx\), on its original
source domain.  Choose a nonempty finite packet of actual zeros and retain
the full order \(m_\rho\) at every selected centre:
\begin{equation}
\begin{gathered}
h(t)=\prod_{\rho\in Z}(t-\rho)^{m_\rho},\qquad d=\deg h,\\
E_h=\mathbb C[t]/(h),\qquad\upsilon_h=[g/h]_h\in E_h^\times.
\end{gathered}
\tag{FD.2}\label{FD:FD:packet}
\end{equation}
The bracket \([\cdot]_h\) means the complete local Taylor jets modulo
\(h\); it includes every nilpotent jet at a multiple zero.  Since each
selected order is the full order of that zero, \(g/h\) is nonzero at every
selected centre, proving that \(\upsilon_h\) is a unit in the finite
quotient.  Let \(A_h\) be multiplication by \([t]_h\) on \(E_h\).

The retained actual source inverse supplies \(F_h\) with
\begin{equation}
h(D)F_h=f_0,\qquad \mathcal MF_h=g/h,\qquad
\mathcal T_h(P)=P(D)F_h.
\tag{FD.3}\label{FD:FD:source}
\end{equation}
Its exact identities, retained from the original source construction, are
\begin{equation}
\mathcal M\mathcal T_h(P)=(g/h)P,\qquad
\mathcal T_h(hP)=\Theta(P(D)\phi_*),\qquad
J_h\mathcal T_h(P)=\upsilon_h[P]_h.
\tag{FD.4}\label{FD:FD:source-identities}
\end{equation}
Here \(J_hF=[\mathcal MF]_h\) is the full Mellin-jet map on the original
analytic source.  The second identity gives an actual theta primitive, and the last identity
includes the arithmetic coefficient \(\upsilon_h\).  The map
\(\mathcal T_h\) is linear and injective: a zero image has identically
zero Mellin transform \((g/h)P\), which forces \(P=0\).

The Hilbert observation remains
\(\mathcal H=L^2((0,\infty),dx)\).  The exact change of variables and
Fourier map
\[
F\longmapsto \bigl(y\longmapsto e^{y/2}F(e^y)\bigr)
\longmapsto\left(t\longmapsto\int_{\mathbb R}e^{y/2}F(e^y)e^{ity}\,dy\right)
\]
have target measure \(dt/(2\pi)\) and composite
\(F\mapsto\mathcal MF(1/2+it)\).  Consequently the arithmetic polynomial
measure is exactly
\begin{equation}
d\nu_h(y)=\left|\frac{g(\tfrac12+iy)}{h(\tfrac12+iy)}\right|^2
\frac{dy}{2\pi}.
\tag{FD.5}\label{FD:FD:measure}
\end{equation}
At selected zeros on the integration line, the quotient is evaluated as
the entire function \(g/h\).  The finiteness of every polynomial moment
can also be checked directly from the retained factors.  If
\(\eta(s)=\sum_{n\ge1}(-1)^{n-1}n^{-s}\), its alternating coefficient
partial sums have absolute value at most one.  Partial summation gives
\(\lvert\eta(s)\rvert\le\lvert s\rvert/\Re s\) for \(\Re s>0\).
The identity
\(\eta(s)=(1-2^{1-s})\zeta(s)\), first obtained by splitting the
absolutely convergent series and then continued analytically, implies
\[
|\zeta(\tfrac12+iy)|
\le\frac{2|\tfrac12+iy|}{\sqrt2-1},
\]
because \(|1-2^{1/2-iy}|\ge\sqrt2-1\).
Fix \(0<\theta<\pi/2\).  Rotating the Euler gamma integral through
angle \(\operatorname{sgn}(y)\theta\) gives
\[
|\Gamma(\tfrac14+iy/2)|
\le\Gamma(\tfrac14)(\cos\theta)^{-1/4}e^{-\theta|y|/2}.
\]
For \(y=0\), the same bound follows from the unrotated integral.
The rotation follows by integrating
\(z^{1/4+iy/2-1}e^{-z}\) around a truncated sector using its continuous
logarithm there.  The small circular arc tends to zero as its radius
to the power \(1/4\); the large arc tends to zero because
\(\cos\theta>0\).  Taking absolute values on the rotated ray yields
the displayed bound.  Retaining the factors in \eqref{FD:FD:arithmetic}
therefore gives
\[
|g(\tfrac12+iy)|^2
\le
\frac{4\pi^{-1/2}\Gamma(\tfrac14)^2}
{(\sqrt2-1)^2\sqrt{\cos\theta}}\,
(\tfrac14+y^2)^3e^{-\theta|y|}.
\]
Monicity of the fixed degree-\(d\) polynomial \(h\) gives
\(|h(\tfrac12+iy)|\ge c_h(1+|y|)^d\) for some \(c_h>0\) outside
a sufficiently large interval: divide its value by
\((\tfrac12+iy)^d\), whose quotient tends to one.  On the remaining
compact interval, the entire function \(g/h\) is bounded.
These estimates prove finiteness of every moment of \(\nu_h\).
The measure is positive away from a discrete set.
Therefore the norm of a nonzero polynomial is
strictly positive: zero norm would make the polynomial vanish almost
everywhere on a line and hence identically.

Fix \(k\ge1\), and let
\begin{equation}
\begin{gathered}
E^{(k)}=E_h^{\otimes k},\qquad
A_k=\sum_{j=1}^k A_{h,j},\qquad D^{(k)}=\sum_{j=1}^kD_j,\\
\mathcal P_M^{(k)}=\mathbb C[t_1,\ldots,t_k]_{\deg\le M},\qquad
\mathcal T^{(k)}(P)=P(D_1,\ldots,D_k)F_h^{\otimes k},\\
\mathcal J(P)=\upsilon_h^{\otimes k}
[P]_{(h(t_1),\ldots,h(t_k))}.
\end{gathered}
\tag{FD.6}\label{FD:FD:tensor-source}
\end{equation}
All tensor generators in this section are additive sums, as displayed in
\eqref{FD:FD:tensor-source}.  The multiplicative symmetric power of a linear
map is a different construction and is not used here.

For \(M\ge k(d-1)\), the map \(\mathcal J\) on
\(\mathcal P_M^{(k)}\) is onto: each reduced monomial has each coordinate
degree below \(d\), hence total degree at most \(k(d-1)\), and multiplication
by \(\upsilon_h^{\otimes k}\) is invertible.  Its kernel and its actual
theta image are
\begin{equation}
\mathcal I_M=\mathcal P_M^{(k)}\cap(h(t_1),\ldots,h(t_k)),\qquad
\mathcal B_M=\mathcal T^{(k)}\mathcal I_M.
\tag{FD.7}\label{FD:FD:relations}
\end{equation}
Successive monic division in the variables expresses every polynomial in
the kernel as \(\sum_jh(t_j)P_j\).  Under \(\mathcal T^{(k)}\), each
summand is an actual theta boundary in factor \(j\), by
\eqref{FD:FD:source-identities}.  Its cochain primitive has coefficient
\((-1)^{j-1}\); the tensor differential has the same coefficient, so their
product gives the displayed positive summand \(+h(t_j)P_j\).

Let \(R_M:E^{(k)}\to\mathcal T^{(k)}\mathcal P_M^{(k)}\) be the unique
least-norm lift of the full jet, using the product measure induced by
\eqref{FD:FD:measure}, and write
\(G_M=R_M^*R_M\), \(\mathcal C_M=G_M^{-1}\).
The existence, uniqueness, and inverse are also obtained directly below.
The minimum is orthogonal to the specified \(\mathcal B_M\), and its
full jet is identity.  Its arithmetic class is that of the original
product source lift with the specified jet, since all admissible
changes are the actual boundaries in \eqref{FD:FD:relations}.

Start with the monic, unscaled orthogonal polynomials and their actual
norms
\begin{equation}
p_0=1,\quad
p_n=t^n-\operatorname{proj}_{\mathbb C[t]_{\le n-1}}t^n,
\quad \omega_n=\int|p_n(\tfrac12+iy)|^2\,d\nu_h(y)>0.
\tag{FD.8}\label{FD:FD:polynomials}
\end{equation}
For a multi-index \(\alpha\in\mathbb N_0^k\), set
\begin{equation}
p_\alpha=\prod_{j=1}^kp_{\alpha_j}(t_j),\qquad
\omega_\alpha=\prod_{j=1}^k\omega_{\alpha_j},\qquad
z_\alpha=\mathcal J(p_\alpha).
\tag{FD.9}\label{FD:FD:product-jets}
\end{equation}
Let \(T_M\) have columns \(\mathcal T^{(k)}p_\alpha\) for
\(|\alpha|\le M\), let \(J_M=[z_\alpha]_{|\alpha|\le M}\), and let
\(\Omega_{\le M}=\operatorname{diag}_{|\alpha|\le M}\omega_\alpha\).
Product orthogonality gives \(T_M^*T_M=\Omega_{\le M}\).
Surjectivity of \(J_M\) implies
\begin{equation}
\begin{gathered}
\mathcal C_M=J_M\Omega_{\le M}^{-1}J_M^*
=\sum_{|\alpha|\le M}\frac{z_\alpha z_\alpha^*}{\omega_\alpha}>0,\\
G_M=\mathcal C_M^{-1},\qquad
R_M=T_M\Omega_{\le M}^{-1}J_M^*G_M.
\end{gathered}
\tag{FD.10}\label{FD:FD:minimum}
\end{equation}
Indeed \(J_M\Omega_{\le M}^{-1}J_M^*G_M=I\).  Any other coefficient
vector with jet \(u\) is \(\Omega_{\le M}^{-1}J_M^*G_Mu+w\), where
\(J_Mw=0\).  Its cross term with the first vector is
\(w^*J_M^*G_Mu=0\), so its squared norm is exactly
\(u^*G_Mu+w^*\Omega_{\le M}w\).  This proves the unique minimum,
its orthogonality to \(\mathcal B_M\), and \(R_M^*R_M=G_M\).

\subsection{The actual next relation and its recurrence bound}

For \(|\beta|=M+1\), define the frontier matrices
\begin{equation}
\Omega=\operatorname{diag}_{|\beta|=M+1}\omega_\beta,\qquad
F=[z_\beta]_{|\beta|=M+1},\qquad
T_+=[\mathcal T^{(k)}p_\beta]_{|\beta|=M+1}.
\tag{FD.11}\label{FD:FD:frontier}
\end{equation}
There are \(r=\binom{M+k}{k-1}\) columns.  They are orthogonal to the
entire degree-\(M\) source image.  If \(P_{\mathcal B_M}\) denotes
orthogonal projection onto the actual boundary space, put
\begin{equation}
\mathcal E_M=(1-P_{\mathcal B_M})\mathcal B_{M+1},\qquad
b_M=T_+-R_MF:\mathbb C^r\longrightarrow\mathcal E_M.
\tag{FD.12}\label{FD:FD:boundary-map}
\end{equation}
The full jet of \(b_M\) is \(F-F=0\), and its two terms are orthogonal
to \(\mathcal B_M\), so its range is contained in \(\mathcal E_M\).
Its highest orthogonal polynomial component is \(T_+\), proving
injectivity.  Conversely, extract the highest component of a member of
\(\mathcal E_M\) and subtract its corresponding \(b_M\)-image.  The
remainder is simultaneously in \(\mathcal B_M\) and orthogonal to
\(\mathcal B_M\), and is therefore zero.  This proves that \(b_M\) is
an isomorphism and constructs its inverse.

The exact Gram and cross-pairing are
\begin{equation}
H_M^{\mathrm{rel}}:=b_M^*b_M=\Omega+F^*G_MF,
\qquad b_M^*R_M=-F^*G_M.
\tag{FD.13}\label{FD:FD:relation-gram}
\end{equation}
The superscript distinguishes this relation Gram from the Hermitian
control matrix introduced later.  The projection of \(R_M\) onto the
new layer and the updated representative are
\begin{equation}
Y_M=-b_M(H_M^{\mathrm{rel}})^{-1}F^*G_M,
\qquad R_{M+1}=R_M-Y_M.
\tag{FD.14}\label{FD:FD:representative-update}
\end{equation}
These corrections are actual boundaries.  In particular, the full jet and
the original arithmetic class stay fixed.  The quotient-layer map
\[
\mathcal B_{M+1}/\mathcal B_M\longrightarrow\mathcal E_M,
\qquad [z]\longmapsto(1-P_{\mathcal B_M})z
\]
is an isomorphism, whose inverse is the inclusion followed by quotient.
Its further map into the quotient by \(\mathcal B_{M+1}\) is zero.
Write \(\lambda_M\) for the specified original joint support label
recording the relation space \(\mathcal B_M\), and write \(0\) for
the zero amplitude in that support fibre.  On the original supported
reconstruction this is
\((\lambda_M,u)\mapsto(\lambda_{M+1},0)\).  For \(u\ne0\) the preceding
class is nonzero.  The final class is the supported zero of that fibre;
it is distinct from \(\tau\), the original external-absence element.
The notation \(e=0^\bullet\) denotes the nonabsent represented scalar
zero in the original scalar reconstruction.  Thus the relation Gram
in \eqref{FD:FD:relation-gram} measures the actual relations becoming
represented zero at the next support.

Orthogonal projection and the finite sum in \eqref{FD:FD:minimum} give
\begin{equation}
\begin{gathered}
G_{M+1}=G_M-G_MF(\Omega+F^*G_MF)^{-1}F^*G_M,\\
\mathcal C_{M+1}-\mathcal C_M=F\Omega^{-1}F^*.
\end{gathered}
\tag{FD.15}\label{FD:FD:full-update}
\end{equation}
The second equality also follows simply by adding the degree-\(M+1\)
terms to \eqref{FD:FD:minimum}; it does not require a limiting argument.

Set \(p_{-1}=0\).  Multiplication by \(t\) on the integration line has adjoint multiplication
by \(1-t\).  Orthogonality and monicity therefore give
\begin{equation}
\begin{gathered}
tp_n=p_{n+1}+b_np_n-a_np_{n-1},\qquad
a_n=\omega_n/\omega_{n-1}\ (n\ge1),\quad a_0=0,\\
b_n=\frac{\int(\tfrac12+iy)|p_n(\tfrac12+iy)|^2d\nu_h(y)}{\omega_n},
\qquad \Re b_n=\tfrac12.
\end{gathered}
\tag{FD.16}\label{FD:FD:recurrence}
\end{equation}
For \(n\ge1\), pairing against a polynomial of degree below \(n-1\)
in the expansion of \(tp_n\)
vanishes by the adjoint identity.  Pairing against \(p_{n-1}\) gives
\(-\omega_n/\omega_{n-1}\), and pairing against \(p_n\) gives the
displayed \(b_n\).  For \(n=0\), monicity gives
\(tp_0=p_1+b_0p_0\), and the term \(a_0p_{-1}=0\) is exactly the
declared zero term.  This proves every coefficient in the recurrence.
The minus sign and the possibly nonzero imaginary part of \(b_n\) are
retained.

Define
\begin{equation}
w_\beta=\sum_{j:\beta_j>0}a_{\beta_j}z_{\beta-e_j},\qquad
E_+=[w_\beta]_{|\beta|=M+1}.
\tag{FD.17}\label{FD:FD:lowering-jets}
\end{equation}
Here \(e_j\) is the standard multi-index vector; it is not the supported
scalar \(e=0^\bullet\).  The full arithmetic jet is a module map, so
\(\mathcal J((\sum_jt_j)P)=A_k\mathcal J(P)\).  Inserting
\eqref{FD:FD:recurrence} into \eqref{FD:FD:minimum} proves
\begin{equation}
\mathcal Z_M:=A_k\mathcal C_M+\mathcal C_MA_k^*-k\mathcal C_M
=F\Omega^{-1}E_+^*+E_+\Omega^{-1}F^*.
\tag{FD.18}\label{FD:FD:full-frontier-identity}
\end{equation}
To verify the cancellation explicitly, for every internal edge
\(\alpha\to\alpha+e_j\), the raising term has coefficient
\(1/\omega_\alpha\), while the lowering term has coefficient
\(-a_{\alpha_j+1}/\omega_{\alpha+e_j}=-1/\omega_\alpha\).
Their adjoints cancel as well.  The diagonal contribution is
\(2\sum_j\Re b_{\alpha_j}=k\).  The only uncancelled edges have
\(|\alpha|=M\) and \(|\alpha+e_j|=M+1\), and their coefficients are
exactly those in the right side of \eqref{FD:FD:full-frontier-identity}.

For completeness, the arithmetic derivative defect has the same actual
relation coordinates.  Set
\[
B_M^{\mathrm{der}}=D^{(k)}R_M-R_MA_k,
\qquad C_M^{\mathrm{rel}}=(1-P_{\mathcal B_M})B_M^{\mathrm{der}}.
\]
Its full jet is zero, and its highest polynomial component is
\(T_+\Omega^{-1}E_+^*G_M\).  The latter assertion follows from the
coefficient of \(p_\beta\) in \((\sum_jt_j)R_M\): it is
\(\sum_{j:\beta_j>0}z_{\beta-e_j}^*G_M/\omega_{\beta-e_j}
=w_\beta^*G_M/\omega_\beta\).  The layer isomorphism
\eqref{FD:FD:boundary-map} therefore proves
\begin{equation}
\begin{gathered}
C_M^{\mathrm{rel}}=b_M\Omega^{-1}E_+^*G_M,\\
\begin{aligned}
W_M&:=A_k^*G_M+G_MA_k-kG_M\\
&=G_M\mathcal Z_MG_M\\
&=-(Y_M^*C_M^{\mathrm{rel}}+(C_M^{\mathrm{rel}})^*Y_M).
\end{aligned}
\end{gathered}
\tag{FD.19}\label{FD:FD:actual-control}
\end{equation}
The last equality follows by substituting
\eqref{FD:FD:relation-gram} and \eqref{FD:FD:representative-update}; in
particular its minus sign is retained.

The recurrence constant is
\begin{equation}
\Gamma_{k,M}=\max_{|\alpha|=M}\sum_{j=1}^k
\bigl(a_{\alpha_j+1}+(k-1)a_{\alpha_j}\bigr)>0.
\tag{FD.20}\label{FD:FD:Gamma}
\end{equation}
Strict positivity follows from \(a_{n+1}>0\) for every \(n\ge0\).
We now prove its complete lowering estimate.

For coefficient calculations only, use the exact coordinate map that
divides the basis vector \(p_\alpha\) by \(\sqrt{\omega_\alpha}\).
Its inverse multiplies by \(\sqrt{\omega_\alpha}\); the original
polynomials, their norms, and the measure are unchanged.
The degree-\((-1)\) coefficient space is by definition \(\{0\}\);
in particular every lowering map from degree zero has that zero
codomain.  This retains the admitted endpoint \(d=1,M=0\).
In these coefficient spaces let \(R_j\) be the raising map
\(e_\alpha\mapsto\sqrt{a_{\alpha_j+1}}e_{\alpha+e_j}\), and let
\(R_j^*\) be its lowering adjoint.  Put
\[
Q_M^{\mathrm{jet}}=[z_\alpha/\sqrt{\omega_\alpha}]_{|\alpha|=M},
\qquad B^{\mathrm{inc}}=(\textstyle\sum_jR_j)^*\big|_{M+1\to M}.
\]
The retained coefficients in \eqref{FD:FD:lowering-jets} give exactly
\(E_+\Omega^{-1/2}=Q_M^{\mathrm{jet}}B^{\mathrm{inc}}\).
Raising and lowering in distinct coordinates commute.  Hence, on the
degree-\(M\) coefficient space,
\[
B^{\mathrm{inc}}(B^{\mathrm{inc}})^*
=\operatorname{diag}_{|\alpha|=M}\sum_j
(a_{\alpha_j+1}-a_{\alpha_j})+L^*L,
\qquad L=\textstyle\sum_j R_j^*\big|_{M\to M-1}.
\]
For every coefficient vector \(c\), Cauchy--Schwarz gives
\(\|Lc\|^2\le k\sum_j\|R_j^*c\|^2\).  Thus
\[
L^*L\preceq k\operatorname{diag}_{|\alpha|=M}\sum_ja_{\alpha_j},
\qquad
B^{\mathrm{inc}}(B^{\mathrm{inc}})^*\preceq\Gamma_{k,M}I.
\]
Multiplication by \(Q_M^{\mathrm{jet}}\) and its adjoint, followed by
inclusion of the positive degree-\(M\) sum in \eqref{FD:FD:minimum}, proves
\begin{equation}
E_+\Omega^{-1}E_+^*\preceq
\Gamma_{k,M}Q_M^{\mathrm{jet}}(Q_M^{\mathrm{jet}})^*
\preceq\Gamma_{k,M}\mathcal C_M.
\tag{FD.21}\label{FD:FD:lowering-bound}
\end{equation}

\subsection{Exact symmetric restriction with every orbit factor}

Let \(P_\pi\) be ordinary factor permutation.  The relevant module is
\[
E_k^{\mathrm{sym}}=(E_h^{\otimes k})^{S_k},\qquad
N=\dim E_k^{\mathrm{sym}}=\binom{k+d-1}{d-1}.
\]
It is the actual top-degree summand selected on the original cochain
tensor product by
\(\mathsf S_k=(1/k!)\sum_{\pi\in S_k}\operatorname{sgn}(\pi)T_\pi\),
where \(T_\pi\) is permutation with its Koszul sign.
In top degree \(T_\pi=\operatorname{sgn}(\pi)P_\pi\), so the two signs
cancel.  The numerator squares to \(k!\) times itself, proving the
idempotent with its literal coefficient \(1/k!\).  In particular,
every repeated tensor \(v^{\otimes k}\) remains in this summand.

Use the quotient basis formed by the unscaled orbit sums of
\([t^{a_1}]\otimes\cdots\otimes[t^{a_k}]\), \(0\le a_i<d\).
Let \(I_k\) be its coefficient inclusion in the full tensor basis.
If \(n_j\) counts factors with exponent \(j\), its orbit size is
\(k!/\prod_{j=0}^{d-1}n_j!\).  The orbit supports are disjoint, so
\begin{equation}
D_k=I_k^*I_k
=\operatorname{diag}\left(\frac{k!}{\prod_{j=0}^{d-1}n_j!}\right),
\qquad L_k=D_k^{-1}I_k^*,\qquad L_kI_k=I.
\tag{FD.22}\label{FD:FD:orbit-coordinates}
\end{equation}
The projector onto invariant coefficient vectors is \(\Pi_k=I_kL_k\).
Permutations preserve the source, the relation ideal, the product
measure, and \(\upsilon_h^{\otimes k}\).  Uniqueness of the minimum
therefore gives equivariance of \(R_M\), and consequently of \(G_M\)
and \(\mathcal C_M\).  The exact restricted matrices are
\begin{equation}
G_M^{\mathrm{sym}}=I_k^*G_MI_k,\qquad
\mathcal C_M^{\mathrm{sym}}=L_k\mathcal C_ML_k^*,\qquad
A_kI_k=I_kA_k^{\mathrm{sym}}.
\tag{FD.23}\label{FD:FD:symmetric-matrices}
\end{equation}
The first two matrices are inverse.  For example,
\[
G_M^{\mathrm{sym}}\mathcal C_M^{\mathrm{sym}}
=I_k^*G_M\Pi_k\mathcal C_ML_k^*
=I_k^*\Pi_kL_k^*=I,
\]
because both \(G_M\) and \(\mathcal C_M\) commute with \(\Pi_k\),
their product is identity, \(\Pi_kL_k^*=L_k^*\), and
\(I_k^*L_k^*=(L_kI_k)^*=I\).

In the source frontier let \(Q_+\) be the coefficient inclusion matrix
whose columns are the \(0/1\) orbit indicators in the ordered basis
\((p_\beta)_{|\beta|=M+1}\).  Thus each column represents the
unscaled orbit sum of those \(p_\beta\).  If \(m_j(\beta)\)
counts the entries of \(\beta\) equal to \(j\), the corresponding
source orbit norm is exactly
\begin{equation}
\omega_{[\beta]}=
\frac{k!}{\prod_{j\ge0}m_j(\beta)!}\prod_{i=1}^k\omega_{\beta_i}.
\tag{FD.24}\label{FD:FD:source-orbit-norm}
\end{equation}
This is a sum of equal squared norms of pairwise orthogonal distinct
product polynomials.  Define the actual restricted frontier matrices
\begin{equation}
\Omega_s=Q_+^*\Omega Q_+,\qquad
F_s=L_kFQ_+,\qquad E_s=L_kE_+Q_+.
\tag{FD.25}\label{FD:FD:symmetric-frontier}
\end{equation}
Thus \(\Omega_s\) is the positive diagonal matrix with entries
\eqref{FD:FD:source-orbit-norm}; none of its orbit factors are suppressed.

The source invariant projector is
\(\Pi_+=Q_+(Q_+^*Q_+)^{-1}Q_+^*\).  Since \(\Omega\) is constant
on each orbit, it commutes with \(\Pi_+\), and direct calculation in
one orbit gives
\[
Q_+\Omega_s^{-1}Q_+^*=\Omega^{-1}\Pi_+.
\]
Indeed, on an orbit of size \(m\) and weight \(\omega\), the left
side is the all-ones matrix divided by \(m\omega\); this is precisely
\(\omega^{-1}\Pi_+\) there.  Equivariance implies
\(L_kF=L_kF\Pi_+\) and \(L_kE_+=L_kE_+\Pi_+\): averaging any
source vector before applying an equivariant map has the same invariant
component as applying the map first.  Congruence of
\eqref{FD:FD:full-update}, \eqref{FD:FD:full-frontier-identity}, and
\eqref{FD:FD:lowering-bound} therefore gives, with
\begin{equation}
\begin{gathered}
\mathcal C=\mathcal C_M^{\mathrm{sym}},\quad
G=G_M^{\mathrm{sym}}=\mathcal C^{-1},\quad A=A_k^{\mathrm{sym}},\\
\Delta\mathcal C=\mathcal C_{M+1}^{\mathrm{sym}}-
\mathcal C_M^{\mathrm{sym}},\qquad
\mathcal Z=A\mathcal C+\mathcal CA^*-k\mathcal C,
\end{gathered}
\tag{FD.26}\label{FD:FD:abbreviations}
\end{equation}
the exact formulas
\begin{equation}
\begin{gathered}
\Delta\mathcal C=F_s\Omega_s^{-1}F_s^*,\\
\mathcal Z=F_s\Omega_s^{-1}E_s^*+E_s\Omega_s^{-1}F_s^*,\\
E_s\Omega_s^{-1}E_s^*\preceq\Gamma_{k,M}\mathcal C.
\end{gathered}
\tag{FD.27}\label{FD:FD:restricted-identities}
\end{equation}
For the operator term, equivariance also gives
\(L_kA_k=A_k^{\mathrm{sym}}L_k\), proving that the congruence of
\(A_k\mathcal C_M+\mathcal C_MA_k^*-k\mathcal C_M\) is exactly the
\(\mathcal Z\) in \eqref{FD:FD:abbreviations}.

\subsection{The aggregate spectral-position restriction}

Positive square roots below are the unique positive square roots of the
retained positive matrices.  They are used as explicit comparison maps;
their inverses are retained, and neither the source measure nor the
arithmetic metric is redefined.  Set
\begin{equation}
\begin{gathered}
X=G^{1/2}F_s\Omega_s^{-1/2},\qquad
Y=G^{1/2}E_s\Omega_s^{-1/2},\\
K=XX^*=G^{1/2}\Delta\mathcal C G^{1/2},\qquad
H=G^{1/2}\mathcal ZG^{1/2}=XY^*+YX^*.
\end{gathered}
\tag{FD.28}\label{FD:FD:comparison-maps}
\end{equation}
The last inequality of \eqref{FD:FD:restricted-identities} gives
\(YY^*\preceq\Gamma_{k,M}I\), and hence
\(Y^*Y\preceq\Gamma_{k,M}I\) on its own coefficient space.  Here the
latter statement follows from \(\|Y\|_{\mathrm{op}}^2
=\|YY^*\|_{\mathrm{op}}\), or directly from the singular values.
For the Frobenius norm \(\|T\|_F^2=\operatorname{tr}(T^*T)\),
\[
\|XY^*\|_F^2
=\operatorname{tr}(X^*X\,Y^*Y)
\le\Gamma_{k,M}\operatorname{tr}(X^*X).
\]
The norm of a matrix equals the norm of its adjoint, so the triangle
inequality gives the proved upper estimate
\begin{equation}
\|H\|_F^2\le4\Gamma_{k,M}\operatorname{tr}K.
\tag{FD.29}\label{FD:FD:Frobenius-upper}
\end{equation}

List the roots of \(h\), repeated with their original algebraic
multiplicities, as \(\rho_1,\ldots,\rho_d\), and put
\begin{equation}
\delta_i=\Re\rho_i-\tfrac12,\qquad
\eta_{\mathbf n}=\sum_{i=1}^dn_i\rho_i,
\quad \mathbf n\in\mathbb N_0^d,\quad |\mathbf n|=k.
\tag{FD.30}\label{FD:FD:spectral-positions}
\end{equation}
The eigenvalues of \(A\), with algebraic multiplicity, are exactly the
\(\eta_{\mathbf n}\).  To prove this while retaining nilpotents,
triangularize \(A_h\) by induction using an eigenvector and the induced
operator on the quotient.  In a basis with
\(A_he_j=\rho_je_j+\sum_{i<j}c_{ij}e_i\), applying the additive
generator to a symmetric monomial of occupations \(\mathbf n\) gives
diagonal coefficient \(\sum_jn_j\rho_j\).  Every other term replaces
an index \(j\) by \(i<j\), reducing the sum of the indices.  Ordering
monomials by that sum, and resolving ties arbitrarily, makes the
operator triangular with precisely the displayed diagonal.  Each weak
composition contributes once; equal eigenvalues remain repeated with
their algebraic multiplicities.  The triangular off-diagonal entries
are retained, including those caused by nontrivial Jordan blocks.

Let \(B=G^{1/2}AG^{-1/2}\).  Substitution into
\eqref{FD:FD:abbreviations} gives \(H=B+B^*-kI\).
Unitary Schur triangularization gives \(T=U^*BU\), upper triangular
with diagonal entries \(\eta_{\mathbf n}\).  Such a triangularization
is obtained inductively by taking a unit eigenvector as the first basis
vector and applying the induction to the compression on its orthogonal
complement.  Unitary conjugation preserves the Frobenius norm.  Its
diagonal and off-diagonal entries give the exact identity
\begin{equation}
\|H\|_F^2
=\sum_{|\mathbf n|=k}(2\Re\eta_{\mathbf n}-k)^2
+2\sum_{i<j}|T_{ij}|^2
\ge4\sum_{|\mathbf n|=k}\left(\sum_i n_i\delta_i\right)^2.
\tag{FD.31}\label{FD:FD:Schur-exact}
\end{equation}
For \(i<j\), the two corresponding entries of \(T+T^*-kI\) are
\(T_{ij}\) and \(\overline{T_{ij}}\); this proves the factor two.
Thus the retained nilpotent and nonnormal contributions cannot weaken
the displayed lower bound.

We evaluate its sum exactly.  For \(d\ge2\), the ordinary generating
functions for weak compositions give
\begin{equation}
\begin{aligned}
\sum_{|\mathbf n|=k}n_1&=\binom{k+d-1}{d},\\
\sum_{|\mathbf n|=k}n_1(n_1-1)&=2\binom{k+d-1}{d+1},\\
\sum_{|\mathbf n|=k}n_1n_2&=\binom{k+d-1}{d+1}.
\end{aligned}
\tag{FD.32}\label{FD:FD:composition-moments}
\end{equation}
Indeed their generating functions are respectively
\[
z(1-z)^{-d-1},\qquad
2z^2(1-z)^{-d-2},\qquad
z^2(1-z)^{-d-2}.
\]
These follow by differentiating the appropriate geometric-series
factors.  Extracting the coefficient of \(z^k\)
gives \eqref{FD:FD:composition-moments}.  Write
\(a=\binom{k+d-1}{d}\) and \(b=\binom{k+d-1}{d+1}\).
Then the square expands as
\[
(a+2b)\sum_i\delta_i^2+2b\sum_{i<j}\delta_i\delta_j
=(a+b)\sum_i\delta_i^2+b\left(\sum_i\delta_i\right)^2.
\]
Pascal's identity therefore proves
\begin{equation}
\begin{aligned}
\mathcal S_{h,k}
&:=\sum_{|\mathbf n|=k}\left(\sum_i n_i\delta_i\right)^2\\
&=\binom{k+d}{d+1}\sum_i\delta_i^2
+\binom{k+d-1}{d+1}\left(\sum_i\delta_i\right)^2.
\end{aligned}
\tag{FD.33}\label{FD:FD:composition-sum}
\end{equation}
Binomial coefficients with smaller upper than nonnegative lower index
are zero.  For \(d=1\), direct evaluation gives \(k^2\delta_1^2\);
the right side of \eqref{FD:FD:composition-sum} is
\([\binom{k+1}{2}+\binom{k}{2}]\delta_1^2=k^2\delta_1^2\), so the
same formula holds in that case.

Combining \eqref{FD:FD:Frobenius-upper}, \eqref{FD:FD:Schur-exact}, and
\eqref{FD:FD:composition-sum} proves the aggregate arithmetic restriction
\begin{equation}
\boxed{\Gamma_{k,M}\operatorname{tr}\!\left(
G_M^{\mathrm{sym}}\bigl(\mathcal C_{M+1}^{\mathrm{sym}}
-\mathcal C_M^{\mathrm{sym}}\bigr)\right)\ge\mathcal S_{h,k}.}
\tag{FD.34}\label{FD:FD:aggregate-bound}
\end{equation}
The trace in this formula is exactly \(\operatorname{tr}K\), by
cyclicity.  All its matrices were defined from the actual measure,
complete arithmetic jet unit, and unscaled orbit norms.

\subsection{Coordinate invariance and determinant growth}

The comparison maps have an exact coordinate-invariance statement.
For invertible changes \(S\) in quotient coordinates and \(R\) in
frontier coordinates, set
\begin{equation}
\begin{gathered}
A'=SAS^{-1},\quad \mathcal C'=S\mathcal CS^*,\quad
G'=S^{-*}GS^{-1},\\
F_s'=SF_sR,\quad E_s'=SE_sR,\quad
\Omega_s'=R^*\Omega_sR.
\end{gathered}
\tag{FD.35}\label{FD:FD:coordinate-change}
\end{equation}
Here \(S^{-*}=(S^{-1})^*\).  Cancellation of \(R,R^{-1}\) in the
factorizations proves \(\Delta\mathcal C'=S\Delta\mathcal CS^*\),
and substitution in the generator identity gives
\(\mathcal Z'=S\mathcal ZS^*\).  The lowering inequality is carried
through the same congruence.  Direct cyclic cancellation and determinant
multiplicativity give
\begin{equation}
\begin{gathered}
\operatorname{tr}(G'\Delta\mathcal C')
=\operatorname{tr}(G\Delta\mathcal C),\\
\frac{\det(\mathcal C'+\Delta\mathcal C')}{\det\mathcal C'}
=\frac{\det(\mathcal C+\Delta\mathcal C)}{\det\mathcal C}.
\end{gathered}
\tag{FD.36}\label{FD:FD:coordinate-invariants}
\end{equation}
Furthermore
\(U_S=(\mathcal C')^{-1/2}S\mathcal C^{1/2}\) satisfies
\(U_SU_S^*=I\), and hence is unitary, with
\(H'=U_SHU_S^*\).  This follows by substituting
\(\mathcal C'=S\mathcal CS^*\) and
\(\mathcal Z'=S\mathcal ZS^*\).
Thus the exact trace, determinant ratio, and Frobenius quantity are
preserved by the comparison-coordinate maps; no arithmetic
normalization has been imposed.

The exact update and \eqref{FD:FD:comparison-maps} give
\begin{equation}
\frac{\det\mathcal C_{M+1}^{\mathrm{sym}}}
{\det\mathcal C_M^{\mathrm{sym}}}=\det(I+K).
\tag{FD.37}\label{FD:FD:determinant-identity}
\end{equation}
Indeed congruence by \(\mathcal C^{-1/2}=G^{1/2}\) takes
\(\mathcal C+\Delta\mathcal C\) to \(I+K\), and its determinant
multiplier is \(1/\det\mathcal C\).  If
\(\kappa_1,\ldots,\kappa_N\ge0\) are the eigenvalues of \(K\),
expanding \(\prod_j(1+\kappa_j)\) gives
\(\det(I+K)\ge1+\sum_j\kappa_j=1+\operatorname{tr}K\).
Since \(\Gamma_{k,M}>0\) was proved in \eqref{FD:FD:Gamma},
\eqref{FD:FD:aggregate-bound} now proves
\begin{equation}
\boxed{\frac{\det\mathcal C_{M+1}^{\mathrm{sym}}}
{\det\mathcal C_M^{\mathrm{sym}}}
\ge1+\frac{\mathcal S_{h,k}}{\Gamma_{k,M}}.}
\tag{FD.38}\label{FD:FD:determinant-step}
\end{equation}
For any integers \(M_1>M_0\ge k(d-1)\), multiplying the finitely many
positive inequalities \eqref{FD:FD:determinant-step} telescopes to
\begin{equation}
\boxed{\frac{\det\mathcal C_{M_1}^{\mathrm{sym}}}
{\det\mathcal C_{M_0}^{\mathrm{sym}}}
\ge\prod_{M=M_0}^{M_1-1}
\left(1+\frac{\mathcal S_{h,k}}{\Gamma_{k,M}}\right).}
\tag{FD.39}\label{FD:FD:determinant-product}
\end{equation}
This is a finite product over the specified actual support degrees.
It requires no limit or unproved asymptotic estimate.

For a reflection-stable packet, the original reflection maps each zero
\(\rho\) to \(1-\overline\rho\), retaining its full order.  Each
paired displacement \(\delta_i\) is sent to \(-\delta_i\), so
\(\sum_i\delta_i=0\).  Hence the exact coefficient becomes
\begin{equation}
\boxed{\mathcal S_{h,k}=\binom{k+d}{d+1}
\sum_{i=1}^d(\Re\rho_i-\tfrac12)^2.}
\tag{FD.40}\label{FD:FD:reflection-coefficient}
\end{equation}
In particular, the presence of any actual off-axis zero makes this
coefficient strictly positive.  It forces a positive relative frontier
increment and the determinant lower bound at every admitted support
degree, within the actual symmetric summand.

\subsection{The exact signed eigenline witness}

Let \(v\ne0\) be an actual eigenvector of \(A_h\), with
\(A_hv=\rho v\).  In the unscaled symmetric quotient coordinates put
\(u=L_k(v^{\otimes k})\), so \(I_ku=v^{\otimes k}\).
This vector is nonzero: choose a linear functional \(\ell\) with
\(\ell v=1\); then \(\ell^{\otimes k}(v^{\otimes k})=1\).
The additive generator satisfies \(Au=k\rho u\).  With the exact
positive number \(q=u^*Gu\), put
\begin{equation}
x=\Omega_s^{-1/2}F_s^*Gu,\qquad
y=\Omega_s^{-1/2}E_s^*Gu.
\tag{FD.41}\label{FD:FD:witness-frontier-vectors}
\end{equation}
The symmetric restriction of \eqref{FD:FD:actual-control} is
\(W=A^*G+GA-kG=G\mathcal ZG\).  On the eigenvector,
\[
u^*Wu=k(2\Re\rho-1)q.
\]
On the other hand, \eqref{FD:FD:restricted-identities} and
\eqref{FD:FD:witness-frontier-vectors} give
\(u^*G\mathcal ZGu=x^*y+y^*x=2\Re(x^*y)\).
Equating these two exact expressions proves the signed identity
\begin{equation}
\boxed{\Re(x^*y)=k(\Re\rho-\tfrac12)q.}
\tag{FD.42}\label{FD:FD:signed-witness}
\end{equation}
Thus this correlation has the retained sign of the actual real-part
displacement.  No absolute-value estimate was used in obtaining it.

Define the two directional costs and the actual symmetric maximum cost
\begin{equation}
\alpha=\frac{\|x\|^2}{q},\qquad
\beta=\frac{\|y\|^2}{q},\qquad
\lambda_{k,M}^{\mathrm{sym}}
=\max_{w\ne0}\frac{w^*\Delta\mathcal Cw}{w^*\mathcal Cw}
=\lambda_{\max}(K).
\tag{FD.43}\label{FD:FD:relative-costs}
\end{equation}
The last equality follows by the invertible substitution
\(w=\mathcal C^{-1/2}z\) in the Rayleigh quotient.  The lowering
estimate in \eqref{FD:FD:restricted-identities}, evaluated at \(Gu\),
gives \(\beta\le\Gamma_{k,M}\), because
\((Gu)^*\mathcal C(Gu)=q\).  The same substitution into the
maximum defining \(\lambda_{k,M}^{\mathrm{sym}}\) gives
\(\alpha\le\lambda_{k,M}^{\mathrm{sym}}\).  Cauchy--Schwarz applied
to \eqref{FD:FD:signed-witness} gives
\[
\alpha\beta\ge k^2(\Re\rho-\tfrac12)^2.
\]
Consequently the actual arithmetic costs satisfy
\begin{equation}
\boxed{\Gamma_{k,M}\lambda_{k,M}^{\mathrm{sym}}
\ge k^2(\Re\rho-\tfrac12)^2.}
\tag{FD.44}\label{FD:FD:quadratic-cost}
\end{equation}
Every term in the signed identity and this bound is attached to the
actual next original relation through \eqref{FD:FD:boundary-map},
\eqref{FD:FD:actual-control}, and the exact symmetric restrictions.

\paragraph{Exact scope of the resulting restriction.}
Equations \eqref{FD:FD:aggregate-bound}--\eqref{FD:FD:reflection-coefficient}
and \eqref{FD:FD:signed-witness}--\eqref{FD:FD:quadratic-cost} prove the stated
frontier and determinant restrictions for the original arithmetic
family, including all multiplicities and nilpotent actions.  The source
explicitly records the unestimated quantity: A1829, line 5580, states
that no uniform sublinear excess was established; A1859, line 6013,
states that the finite bound's constants were not bounded sublinearly in
tensor degree; A1859, line 6327, states that the uniform sublinear
estimate for the actual arithmetic family was not proved.  No upper
estimate incompatible with \eqref{FD:FD:determinant-step} or
\eqref{FD:FD:quadratic-cost} is asserted here.  These are proved necessary
restrictions on an actual RH counterexample, not an assertion that their
arithmetic values cannot satisfy them.
